The following are the inversion lemmas and helper lemmas used in the proof of Theorem~\ref{thm:flowsto_event_soundness} above,
showing the correspondence properties between the trace based semantics and the program analysis results.
%
\input{appendices/thm-adaptivity/lem_basicinversion}
%
\begin{lem}[While Loop Inversion]
	\label{lem:inv_while}
	For every $\trace, \trace' \in \ftdom, c, c_1, c_2 \in \cdom$ 
	if $ \config{c, \trace} \rightarrow^* \config{c_1; c_2, \trace'}$ and 
	$c_1 \in_c c_2$, 
	then there must exist a $\ewhile$ command in $c_2$ and $c_1$ must shows up in the body of that $\ewhile$ command,
	 i.e., $\exists l \in \mathbb{N}, b \in \mathcal{B}, c_w \in \cdom \sthat  
	(\ewhile [b]^l \edo (c_w)) \in_c c_2 \land c_1 \in_c c_w$.
	%
	\[
	\begin{array}{l}
	\forall \trace, \trace' \in \ftdom, c, c_1, c_2 \in \cdom \sthat 
		\\ \quad
		\config{c, \trace} \rightarrow^* \config{c_1; c_2, \trace'}
		\implies
		\\ \qquad
		(c_1 \in_c c_2
		\land
		\exists l \in \mathbb{N}, b \in \mathcal{B}, c_w \in \cdom \sthat  
		(\ewhile [b]^l \edo (c_w)) \in_c c_2 \land c_1 \in_c c_w)
	\end{array}
	\]
	\end{lem}	
	Proof Summary: trivially by induction on $c$ and enumerate all operational semantic rules.
\begin{proof}
	Take arbitrary $\trace \in \ftdom$, by induction on $c$, we have following cases:
		\caseL{$c = [\assign{x}{\expr}]^l$}
		By the evaluation rule $\rname{assn}$, we have
		$
		{
		\config{[\assign{{x}}{\aexpr}]^{l},  \trace } 
		\xrightarrow{} 
		\config{\clabel{\eskip}^{l'}, \trace \tracecat [({x}, l, v) ]}
		}$.
		\\
		Since there doesn't exist $c_1, c_2 \in \cdom$ satisfying $\eskip = c_1; c_2$, this theorem is vacuously true.
		\caseL{$c = [\assign{x}{\query(\qexpr)}]^l$}
		By the evaluation rule $\rname{query}$, we have
		$
		{
		\config{[\assign{{x}}{\query(\qexpr)}]^{l},  \trace } 
		\xrightarrow{} 
		\config{\clabel{\eskip}^{l'}, \trace \tracecat [({x}, l, \qval, v) ]}
		}$.
		\\
		Since there doesn't exist $c_1, c_2 \in \cdom$ satisfying $\eskip = c_1; c_2$, this theorem is vacuously true.
		\caseL{$c = \eif([b]^{l}, c_1, c_2)$}
		By the evaluation rule $\rname{query}$ and $\rname{if-f}$, and the label consistency, we know:
		\\
		for all possible $c_{t1}$ and $c_{t2}$ 
		such that $c_t$ has the form $c_t = c_{t1};c_{t2}$;
		\\
		all possible $c_{f1}$ and $c_{f2}$ 
		such that $c_f$ has the form $c_f = c_{f1};c_{f2}$,
		\\
		$c_{t1} \notin c_{t1}$ and $c_{f1} \notin c_{f2}$.
		\\
		Then this theorem is vacuously true.
		\caseL{$c = c_{s1};c_{s2}$}
		By label consistency, we know for every $c_1' \in_c c_{s1}$, $c_1' \notin c_{s2}$.
		\\
		Then by the induction hypothesis on $c_{s1}$ and $c_{s2}$ separately, we have this case proved.
		\caseL{$\ewhile [b]^{l} \edo c$}
		By rule $\rname{while-t}$, we have:
		\[
			\config{{\ewhile [b]^{l} \edo (c_w), \trace}}
			\xrightarrow{} 
			\config{{
			c_w; \ewhile [b]^{l} \edo (c_w),
			\trace \tracecat [\event]}}
		\]
		%
		If $c_w$ is a sequence command,
		let $c_1 = c_{w1}$ be the any possible command in this sequence, for all possible $c_{w1}$ and $c_{w2}$ 
		such that $c_w$ has the form $c_w = c_{w1};c_{w2}$.
		\\
		Then we have $c_2 = c_{w2};\ewhile [b]^{l} \edo (c_w)$ and $c_1 \in_c c_2$.
		\\
		And we also have the existence of $l = l_b, b$ and $c_w$, and $\ewhile [b]^{l} \edo (c_w) \in_c c_2$ and  $c_1 \in c_w$.
		\\
		If $c_w$ isn't a sequence command, let $c_1 = c_w$, then we have $c_2 = \ewhile [b]^{l} \edo (c_w)$ 
		and $c_1 \in_c c_2$.
		\\
		And we also have the existence of $l = l_b, b$ and $c_w$, and $\ewhile [b]^{l} \edo (c_w) \in_c c_2$ and  $c_1 \in c_w$.
		\\
		This case is proved.
		\\
		By the evaluation rule $\rname{while-f}$, we have
		$$
		{
			\config{{\ewhile [b]^{l} \edo (c_w), \trace}}
			\xrightarrow{} 
			\config{{
			\clabel{\eskip}^l ,
			\trace \tracecat [(b, l, \efalse)]}}.
		}$$
		\\
		Since there doesn't exist $c_1, c_2 \in \cdom$ satisfying $\eskip = c_1; c_2$, this theorem is vacuously true.
	\end{proof}
%
%
% \begin{lem}[Only $\eskip$ Command doesn't Produce Event].
% 	\label{lem:inv_skip}
% 	For all trace $\trace\in \ftdom$, and $c, c' \in \cdom$,  
% 	$\config{c, \trace} \rightarrow \config{c', \trace}$ if and only if $c = \clabel{\eskip}^{l'};c'$. 
% 	\[
% 		\forall \trace\in \ftdom, c, c' \in \cdom \sthat 
% 		\config{c, \trace} \rightarrow \config{c', \trace}
% 		\Leftrightarrow 
% 		c = \clabel{\eskip}^{l'};c'
% 	% \footnote{$(\clabel{\eskip}^{l'};){}^*$ denotes a sequence command only composed of $\clabel{\eskip}^{l'}$ commands.}
% 	\]
% 	\end{lem}
% \begin{proof}
% 	Proved trivially by induction on $c$ and enumerate all operational semantic rules.
% \end{proof}
% %
% % \todo{Event Dependency Transitivity}
% \begin{lem}
% 	\label{lem:valdep_trans}
% (Event Dependency Transitivity)
% For every $D \in \dbdom , c \in \cdom, \trace \in \ftdom$, and $\event_1, \event_2, \event_3 \in \eventset^{\asn}, \trace_{12}, \trace_{23} \in \ftdom$,
% if $\eventdep(\event_1, \event_2, \trace_{12}, c, D)$
% and $\eventdep(\event_2, \event_3, \trace_{23}, c, D) $,
% then there is an event dependency relation,
% $\eventdep(\event_1, \event_3, \trace_{12}\tracecat\trace_{23}, c, D)$.
%   % An event $\event_2 \in \eventset^{\asn}$ is in the \emph{may-dependency} relation with another
%   % event $\event_1 \in \eventset^{\asn}$ in a program ${c}$ with a hidden database $D$, denoted as $\eventdep(\event_1, \event_2, c, D)$,
%   % if and only if
%   \[
% 	  \begin{array}{l}
%   \forall D \in \dbdom , c \in \cdom, \event_1, \event_2, \event_3 \in \eventset^{\asn}, \trace_{12}, \trace_{23} \in \ftdom \sthat  
%   \\ \quad
%   \eventdep(\event_1, \event_2, \trace_{12}, c, D) 
%   \land
%   \eventdep(\event_2, \event_3, \trace_{23}, c, D) 
%   \implies
%   \eventdep(\event_1, \event_3, \trace_{12}\tracecat\trace_{23}, c, D)
% 	  \end{array}
%   \]
% \end{lem}
% %
% %

% %
% \begin{lem}[\emph{Variable May-Dependency} Transitivity]
% 	\label{lem:vardep_trans}
% For every $c \in \cdom, x^i, y^j, z^l \in \lvar(c)$, 
% if $\vardep(x^i, y^j, c)$ and 
% $\vardep(y^j, z^l, c)$, then $\vardep(x^i, z^l, c)$.
% 	\[
% 	\forall c \in \cdom, x^i, y^j, z^l \in \lvar(c) \sthat  
% 	\vardep(x^i, y^j, c) 
% 	\land
% 	\vardep(y^j, z^l, c) 
% 	\implies
% 	\vardep(x^i, z^l, c)
% 	\]
%   \end{lem}
  %
%
